\begin{afterword}
\summaryhead

The post starts from Steve Omohundro's principle that a rational agent should spread its resources so that each use brings the same marginal benefit. Such an agent has a common currency of expected utilons. In our society, the author says, that currency is money, the measure of how much society cares about something.

Many people would rather volunteer time or give old goods than give money. The author grants that spending hurts, then recalls the lawyer who volunteers an hour at a soup kitchen instead of working an extra hour and paying for five hours of a hired worker. Comparative advantage, specialization and gains from trade are how things get done, and money is how we use them. Volunteering makes sense for people who cannot sell more of their time, or when a charity needs a large block of the volunteer's own skill. The lawyer may also volunteer to stay motivated, if money is given too. Giving money is not a lesser form of caring. Money measures relative caring, and someone never willing to spend any of it on a thing does not care about it.

\responsehead

The practical advice is sound, and the post limits it with care. An hour of a high earner's pay buys more of a charity's routine work than an hour of that person's own labour. The post states the two main exceptions itself, people who cannot sell more of their time and a specialist's skill given in a block, and it allows volunteering as a way to stay motivated. The moral case for giving money to distant strangers had been made by Peter Singer in 1972, and the case for earning more in order to give more by Peter Unger in 1996, with a law professor as his example.

The title claim is stated more broadly than the argument supports. Omohundro's rule is about one agent's budget. Applied to a society, money measures caring weighted by how much money each person has, as economists who measure benefits in money acknowledge. The post concedes that society is not sane, and later narrows the claim to relative caring within one person's spending, but the opening sentence, that money ``is the measure of how much society cares about something,'' stands as written.

The claim that specialization and trade are ``the only way that anything ever gets done in this world'' is also broader than its support. Benkler documented large projects built from unpaid contributions of a few minutes each, such as NASA's Clickworkers. His condition, that the work can be split into small pieces, leaves the lawyer example standing, but not the general claim.

The post's diagnoses of other people are asserted. It says many are ``motivated to deny'' its point, calls the view that money donors care less ``a mere reflection of a mind that doesn't understand'' markets, and suggests, with ``may,'' that this one misunderstanding could prevent collective action in groups larger than 40 people. None comes with evidence.

For information, with hindsight: an argument like this one became known in effective altruism as ``earning to give.'' In 2015 80,000 Hours, which had promoted it, wrote that only a small proportion of people should do it long term, because effective altruism organizations reported lacking talent more than money. That is the post's own exception, reported as common.

In short: sound and carefully limited advice to give money rather than scattered hours, along lines Singer and Unger had drawn before, set inside larger claims, that money measures society's caring and that specialization is the only way anything gets done, which are stated more broadly than the argument supports.
\end{afterword}
